\section{The derivative the surface does not carry}
\label{sec:ssrmissing}

Take stock of what the book can now do.  Part~II turned a raw board
into a clean observation: Chapter~6 fitted the forward and the
discount, Chapter~7 removed the early-exercise right, Chapter~8 put
the readings on an honest clock.  Part~I turned such observations into
surfaces: at every quoted expiry, a fitted smile, a full probability
law.  All of it is a photograph.  Every object was built from the
prices of one afternoon, at one level of the underlying.

Look at how much the photograph contains.  Differentiate a fitted
smile's call prices once in strike and you get the digital price ---
the probability of finishing below the strike (Chapter~2 built its
whole model on this).  Differentiate twice and you get the density.
Differentiate total variance in maturity and you get the forward
variance of Chapter~8.  In the surface's own coordinates --- strike
and expiry --- every derivative is available, exact, and priced.
Chapter~8 closed on precisely this point, and it is the one aphorism
this chapter allows itself: the surface gives every derivative in the
strike direction and none in the spot direction.

Here is the question that needs the missing direction.  The frozen SPY
December 2026 smile --- the book's running node, \MacSsrHeroDays\ days
to expiry, ATM volatility \MacSsrHeroAtmPct\% --- was fitted at the
snapshot's forward.  Suppose the next tick takes the forward down
$\MacSsrFanMovePct\%$.  What is the new smile, \emph{now}?  No new
option quotes have printed yet, and a desk cannot wait for a
recalibration to know its risk.  It needs a rule that turns the old
smile plus the move into a new smile.

The uncomfortable fact is that today's prices do not contain that
rule.

\begin{remark}[A snapshot does not identify dynamics]
\label{rem:ssrident}
Every price quoted today is an integral against today's laws: Chapter~2
wrote each vanilla as an expectation under its expiry's distribution,
and the fit determined those distributions.  A \emph{dynamics} is a
different kind of object: an assignment of a whole surface to each
future level of the underlying.  Consistency with today's quotes
constrains that assignment at exactly one point --- today's level,
where it must reproduce today's surface --- and nowhere else.  To see
this concretely, take any family of arbitrage-free surfaces that
passes through today's surface.  Declaring ``after a move of size $H$
the surface is the family member at $H$'' \emph{is} choosing a
dynamics.  Every instrument quoted today is priced identically under
every such declaration, so today's quotes cannot rank them.  The
derivative of the smile along spot is extra structure.  The market data of one afternoon cannot supply
it; the desk must choose it, or derive it from a model --- which is
where this chapter ends up in \cref{sec:ssrhalf}.
\end{remark}

The choice would not matter much if it only changed a chart between
ticks.  It matters because it is a \emph{hedge ratio}.  The spot
sensitivity of an option has two parts: the Black delta --- the
sensitivity at frozen volatility --- plus vega times the very
derivative the surface does not carry (vega, recall, is the price's
sensitivity to volatility).  Two desks holding the same option against
the same fitted surface, but assuming different smile responses, will
compute different deltas and trade different hedges.
\Cref{sec:ssrdelta} prices the disagreement on the frozen node: about
\MacSsrDeltaGapPts\ delta points at an ordinary put strike.  Nothing
about the choice is cosmetic.

\begin{figure}[htbp]
\centering
\paperfig{\textwidth}{fig_ssr_question.pdf}
\caption{The missing derivative (frozen SPY December 2026 node).
(a)~Today's fitted smile (orange, with the prepared mid quotes as
dots) and three candidate smiles after a $-\MacSsrFanMovePct\%$
forward move, drawn with the transport built in
\cref{sec:ssrdial} --- all three consistent with every price in
panel~(a)'s dots.  (b)~The volatility of one fixed strike (today's
$k=-0.10$ put, fitted today at \MacSsrMarkVolPct\%) against the move
size,
under the three rules: at a $-\MacSsrFanMovePct\%$ move the answers
sit \MacSsrMarkSpreadBp\ vol bp apart.}
\label{fig:ssrquestion}
\end{figure}

\Cref{fig:ssrquestion} draws the predicament before any machinery.  In
panel~(a), the orange curve is today's fitted smile with its prepared
mid quotes (``prepared'' in Part~II's sense: the cleaned, forwarded,
Europeanized observation); the three thinner curves are three
candidate tomorrows after a $-\MacSsrFanMovePct\%$ forward move, each
produced by one of the standard rules this chapter will build, and
each labeled with the name a desk would use.  They stack in the order
of the eventual dial --- after this \emph{down} move the
sticky-local-vol tomorrow sits highest, sticky-moneyness lowest,
sticky-strike between, at every strike drawn --- for reasons
\cref{sec:ssrdial} derives.  Nothing in the orange dots prefers one
of the three: they are different completions of the same photograph
(\cref{rem:ssrident}).  Panel~(b) reads the same disagreement at a
single fixed strike --- the put ten percent below the money, whose
fitted volatility today is \MacSsrMarkVolPct\%.  Each rule makes that strike's tomorrow
volatility a straight line in the move size; the lines share only the
point at zero move.  At $-\MacSsrFanMovePct\%$ they span
\MacSsrMarkSpreadBp\ vol bp (a vol bp, as throughout the book, is
$10^{-4}$ of absolute volatility) --- three and a half volatility
points for the same strike, same surface, same tick.

We now know what is missing and why it must be chosen.  The rest of
the chapter builds the choices properly: first the bookkeeping that
keeps signs honest (\cref{sec:ssrsigns}), then the one-parameter family
of rules and its dial (\cref{sec:ssrdial}), then the price tag in
delta (\cref{sec:ssrdelta}), and then the one construction in this
book that \emph{answers} the question instead of asking it
(\cref{sec:ssrhalf,sec:ssrfield}).
